\chapter{Аудит кандидатов резервуарно-хопфовского интертвинера}
\label{ch:v10-hopf-reservoir-intertwiner-audit}

\section{После выбора перестановки остаются фазы}

Предыдущий гейт показал, что знак $a_c-a_h=\log2$ единственно выбирает
перестановку «горячая: путь один, холодная: путь два». Но полный
интертвинер $U:\mathcal H_{\rm path}\to\mathcal H_{\rm res}$ должен
удовлетворять
\begin{equation}
 Z_{\rm res}U=UZ_{\rm path},\qquad
 Z_{\rm res}=Z_{\rm path}=\operatorname{diag}(-1,1).
\end{equation}
В координатах $(u_{11},u_{12},u_{21},u_{22})$ эта система имеет оператор
\begin{equation}
 C_U=\operatorname{diag}(0,-2,2,0),\qquad
 \rank C_U=2,qquad \operatorname{nullity}C_U=2.
\end{equation}
Следовательно, ориентация обнуляет внедиагональные элементы, но оставляет
два диагональных коэффициента. В вещественном ортогональном секторе
остаются четыре решения $\operatorname{diag}(\pm1,\pm1)$; над комплексными
числами это фазовый тор $U(1)^2$.

\section{Десять кандидатов}

Проверены унаследованный нулевой блок, спектральное сопоставление
проекторов, базисное тождество, относительный знак, KMS-перестановка,
хопфовский incidence-морфизм, модульное сопряжение, ковариация тока,
минимальный portal-блок и целевым образом загруженный родитель.

Шесть критериев требуют правильного типа $\operatorname{Hom}$,
унаследованности, ненулевого отображения, согласия с ориентацией, выбора
фазы и коэффициента и появления в смешанном блоке общего родителя. Матрица
кандидатов имеет полный ранг $6$, но
\begin{equation}
 N_{\rm pass}=0/10,\qquad N_{\rm inherited}=0/10,
 \qquad N_{\rm max}=5/6.
\end{equation}

Спектральные проекторы формально дают тождественное сопоставление уровней,
но не создают межсекторный оператор и не выбирают относительные фазы.
Ближайший кандидат --- член $\frac12\|r-p\|^2$. Его гессиан
\begin{equation}
 H_{\rm target}=\begin{pmatrix}I&-I\\-I&I\end{pmatrix}
\end{equation}
имеет ненулевой смешанный блок и ранг/ядро $2/2$, однако сам интертвинер
уже помещён в этот функционал. Это реализация цели, а не её происхождение.

\begin{theorem}[Ориентация не фиксирует полный интертвинер]
Уравнение $Z_{\rm res}U=UZ_{\rm path}$ имеет двумерное линейное пространство
решений. Поэтому уникальность перестановки не влечёт уникальности фаз и
коэффициентов межсекторного морфизма.
\end{theorem}

\begin{nogocriterion}[Проекторное совпадение не является общим носителем]
Совпадение спектральных проекторов двух изоморфных двумерных пространств
не создаёт элемент $\operatorname{Hom}(\mathcal H_{\rm path},
\mathcal H_{\rm res})$. Для происхождения интертвинера необходим общий
типизированный носитель, на котором смешанный блок является унаследованным
оператором, а не целевой подстановкой.
\end{nogocriterion}

\begin{protocol}\begin{enumerate}
\item проверено кандидатов: \textbf{$10$};
\item ранг критериев: \textbf{$6/6$};
\item максимальная оценка: \textbf{$5/6$};
\item полных проходов: \textbf{$0/10$};
\item происхождение фазы и смешанного блока: \textbf{$0/2$};
\item следующий гейт: \textbf{допуск общего носителя интертвинера}.
\end{enumerate}\end{protocol}

Машинный сертификат сохранён в
\path{s2t_v10_cell_birth_four_volume_induced_newton_breathing_anomaly_hopf_reservoir_intertwiner_candidate_audit_gate_results.json}.